% GENERATED FILE. Edit canonical lessons, then run npm run generate:books.
% canonical-source-hash: fnv1a64:a17aed654b3a639f
% canonical-lessons: BN-W01-na, BN-W01-na-trace, BN-W01-na-guided-copy, BN-W01-na-delayed-copy, BN-W01-na-dictation, BN-W01-aa-matra, BN-W01-aa, BN-W01-ha, BN-W01-ma, BN-W01-sa, BN-W01-ka, BN-W01-hasanta, BN-W01-ra, BN-W01-nomoshkar-read

\chapter{Hanging From the Head-Line, One Piece at a Time}
\label{ch:bn-script-first}
\hlchaptermodality{2}

\begin{chapteropening}
\textbf{By the end of this chapter:} \emph{I can write nine Bengali pieces from memory — six consonants, a vowel sign, its independent twin and the hasanta — and read the greeting I have been saying since the first chapter by assembling it from those pieces, including its conjunct.}

The greeting assembled from seven pieces the reader can each write from memory, including the conjunct and the letter written s but said sh, and introducing nothing new.
\end{chapteropening}

\section[nô]{\bn{ন} — the consonant na}

\label{lesson:BN-W01-na}



\begin{quote}
Nothing is assumed but the words you can already say. This chapter turns them into words you can read.
\end{quote}

\subsection*{Script}
Bengali is an \textbf{abugida}: a consonant already carries a vowel, and you write a vowel mark only when you want a different one.

\begin{quote}
\bn{ন}
\end{quote}

Now the part that catches every reader who comes to Bengali from Hindi or Sanskrit. The inherent vowel here is \textbf{not} \emph{a}. It is \textbf{ô} — the vowel of English \emph{awe}, rounded and open. So this letter alone is \textbf{nô}, not \emph{na}.

Devanagari's bare consonant says \emph{ka}; Bengali's says \emph{kô}. Same family, same system, \textbf{different default}. It is one of the first things that makes Bengali sound like Bengali rather than like Hindi read aloud.

Look at the top of the shape: a \textbf{horizontal bar}, and in running text it joins up across a whole word so the writing hangs from a line — the same head-line Devanagari has, and the one Gujarati erases.

\subsection*{Your turn}
\begin{itemize}
\raggedright
  \item \emph{Write it:} \bn{ন}
  \item \emph{Read it:} every piece this chapter has given you so far, in order
\end{itemize}

\begin{tcolorbox}[breakable,colback=teal!4,colframe=teal!35!black,title={Before you move on}]
What vowel does a bare Bengali consonant carry? (\textbf{ô}, the vowel of \emph{awe} — not \emph{a}.) And what does Devanagari's carry instead? (\emph{\textbf{a}}.)

Source: \href{https://www.unicode.org/charts/PDF/U0980.pdf}{Unicode Bengali chart}.
\end{tcolorbox}
\section[nô]{\bn{ন} — trace it, with the bar last}

\label{lesson:BN-W01-na-trace}



\begin{quote}
Keep \textbf{\bn{ন}} in front of you. It does not get covered in this lesson.
\end{quote}

\subsection*{Writing — observe and trace}
Finger-trace \textbf{\bn{ন}} three times on the page, slowly, then trace it once in the air. Say \emph{nô} after your finger stops each time — \emph{nô}, with the rounded open vowel of English \emph{awe}, not \emph{na}.

Leave the \textbf{bar across the top for last}, every time. The body of the letter first, the horizontal line after it.

\begin{quote}
This book does not tell you where each curve of the body starts or which way it travels. Bengali handwriting is taught with real variation from school to school, and it is not written down here until it can be written down with a source. Tracing what is in front of you needs no such source, and it is how the shape gets into your hand in the meantime.
\end{quote}

The bar is the exception, and it is not a stylistic claim: the \textbf{\bn{মাত্রা}} joins up across a whole word in running text, so it cannot be the first thing drawn on a letter that is about to have neighbours.

\begin{tcolorbox}[breakable,colback=teal!4,colframe=teal!35!black,title={Before you move on}]
Which part did you leave until last? (\textbf{The bar.}) And what does this letter say on its own — \emph{na} or \emph{nô}? (\emph{\textbf{nô}}.)
\end{tcolorbox}
\section[nô]{\bn{ন} — copy it twice, and join the bar}

\label{lesson:BN-W01-na-guided-copy}



\begin{quote}
The model stays visible for all of this lesson. Nothing is hidden today.
\end{quote}

\subsection*{Writing — guided copy}
Copy \textbf{\bn{ন}} once, body first and bar last, as you traced it.

Then copy it \textbf{twice side by side} — and this time draw the bodies first, both of them, and finish with \textbf{one single bar running across the top of both}.

That one continuous line is the whole lesson. In Bengali the bar is not a mark belonging to a letter; it is a line belonging to the \textbf{word}, which is why it is drawn last and why it runs on. A reader's eye follows it to find where one word stops and the next starts, so a hand that gives each letter its own short bar has written two words where it meant one.

\begin{tcolorbox}[breakable,colback=teal!4,colframe=teal!35!black,title={Before you move on}]
Look at your pair. Is there one line across the top, or two? One line is a word. Two is a gap you did not mean to leave.
\end{tcolorbox}
\section[nô]{\bn{ন} — write it with the model covered}

\label{lesson:BN-W01-na-delayed-copy}



\begin{quote}
Look at \textbf{\bn{ন}} once and say \emph{nô}. Then cover it.
\end{quote}

\subsection*{Writing — delayed copy}
With the model covered, wait ten seconds, then write \textbf{\bn{ন}} once from memory — body first, bar last.

Ten seconds is long enough that your hand cannot be tracing a shape it can still see, and short enough that a letter you met a minute ago is still there.

Only after the letter is finished, uncover the model and compare. Repair \textbf{one} thing, the first that differs, and leave the rest.

\begin{tcolorbox}[breakable,colback=teal!4,colframe=teal!35!black,title={Before you move on}]
Check the bar's \textbf{width} before you check the body at all.

A hand writing from memory remembers to put the bar there and draws it only as wide as the body underneath. In running text it has to reach the full width of the letter, because the next letter's bar will meet it — so a short bar is the one mistake that looks completely right on a single letter and falls apart the moment there are two.
\end{tcolorbox}
\section[nô]{\bn{ন} — the letter from the sound}

\label{lesson:BN-W01-na-dictation}



\begin{quote}
Cover every \textbf{\bn{ন}} on the page. Nothing to copy and nothing to uncover.
\end{quote}

\subsection*{Writing — dictation}
Say \emph{nô} aloud, or play it, and write the letter that carries it. Once.

A delayed copy asks your hand to remember a shape. Dictation asks it to turn a \textbf{sound} into a shape, which is what writing is actually for, and the two do not always arrive together.

\begin{tcolorbox}[breakable,colback=teal!4,colframe=teal!35!black,title={Before you move on}]
Uncover the model and compare — then notice which sound you were given.

\textbf{nô}, not \emph{na}. The vowel inside a bare Bengali consonant is the rounded open one, and this is the place that matters: a reader coming from Hindi or Sanskrit hears \emph{na}, writes the same shape, and has learned the letter with the wrong sound attached to it. The shape is shared across the family. The default vowel is not, and dictation is the only stage that can catch the difference, because it is the only one where the sound comes first.
\end{tcolorbox}
\section[-a]{\bn{া} — the vowel sign that replaces the inherent vowel}

\label{lesson:BN-W01-aa-matra}



\begin{quote}
Write the piece from the previous lesson before adding another.
\end{quote}

\subsection*{Script}
To get a different vowel you \textbf{replace} the inherent one with a mark called a \textbf{mātrā}.

\begin{quote}
\bn{া}
\end{quote}

A single upright stroke standing to the \textbf{right} of its consonant, dropping from the head-line. Attached, it swaps the inherent \textbf{ô} for a long open \textbf{a}.

\begin{quote}
\textbf{\bn{ন}} + \textbf{\bn{া}} $\to$ \textbf{\bn{না}}
\end{quote}

\textbf{na} — and that is the word for \textbf{no}, which you have been saying since your first lesson.

One consonant and one mark, and a word you could already say became a word you can \textbf{read}.

\subsection*{Your turn}
\begin{itemize}
\raggedright
  \item \emph{Write it:} \bn{া}
  \item \emph{Read it:} every piece this chapter has given you so far, in order
\end{itemize}

\begin{tcolorbox}[breakable,colback=teal!4,colframe=teal!35!black,title={Before you move on}]
What does a mātrā do to the inherent vowel? (\textbf{Replaces} it.) What word did those two pieces make? (\textbf{\bn{না}} — \emph{na}, \textquotedblleft{}no\textquotedblright{}.)

Source: \href{https://www.unicode.org/charts/PDF/U0980.pdf}{Unicode Bengali chart}.
\end{tcolorbox}
\section[a]{\bn{আ} — the same vowel standing alone}

\label{lesson:BN-W01-aa}



\begin{quote}
Write the piece from the previous lesson before adding another.
\end{quote}

\subsection*{Script}
The mark you learned a moment ago only works \textbf{attached to a consonant}. So what do you write when a word \textbf{begins} with that vowel?

A different character:

\begin{quote}
\bn{আ}
\end{quote}

That is the same \textbf{a} sound, but as an \textbf{independent vowel} — a full letter that stands on its own at the start of a word.

\textbf{These are two different characters, not two styles of one.} A font, a keyboard and this book all treat them as entirely separate, and confusing them is the commonest spelling error in every Indic script.

What decides which one you write is \textbf{position}, not sound:

\begin{itemize}
\raggedright
  \item \textbf{\bn{আ}} begins a word.
  \item \textbf{\bn{া}} rides on a consonant.
\end{itemize}

Two words you already say begin with this letter — the one for \emph{I see you later} and the one for \emph{all right}.

\subsection*{Your turn}
\begin{itemize}
\raggedright
  \item \emph{Write it:} \bn{আ}
  \item \emph{Read it:} every piece this chapter has given you so far, in order
\end{itemize}

\begin{tcolorbox}[breakable,colback=teal!4,colframe=teal!35!black,title={Before you move on}]
Which of the two starts a word? (\textbf{\bn{আ}}.) Which rides on a consonant? (\textbf{\bn{া}}.) Are they the same character? (\textbf{No} — two separate characters, one sound.)

Source: \href{https://www.unicode.org/charts/PDF/U0980.pdf}{Unicode Bengali chart}.
\end{tcolorbox}
\section[hô]{\bn{হ} — the consonant ha}

\label{lesson:BN-W01-ha}



\begin{quote}
Write the piece from the previous lesson before adding another.
\end{quote}

\subsection*{Script}
\begin{quote}
\bn{হ}
\end{quote}

\textbf{hô} — the \emph{h} of English \emph{hat}, plus the inherent vowel you now know is \textbf{ô} rather than \emph{a}.

No new machinery: every rule you hold applies unchanged. \textbf{New shapes are coming; new rules are not.}

You have been saying this letter since your first lesson — it opens the word for \emph{yes}.

\subsection*{Your turn}
\begin{itemize}
\raggedright
  \item \emph{Write it:} \bn{হ}
  \item \emph{Read it:} every piece this chapter has given you so far, in order
\end{itemize}

\begin{tcolorbox}[breakable,colback=teal!4,colframe=teal!35!black,title={Before you move on}]
What is the inherent vowel in \textbf{\bn{হ}}? (\textbf{ô}.) How many new rules did this lesson add? (\textbf{None}.)

Source: \href{https://www.unicode.org/charts/PDF/U0980.pdf}{Unicode Bengali chart}.
\end{tcolorbox}
\section[mô]{\bn{ম} — the consonant ma}

\label{lesson:BN-W01-ma}



\begin{quote}
Write the piece from the previous lesson before adding another.
\end{quote}

\subsection*{Script}
\begin{quote}
\bn{ম}
\end{quote}

\textbf{mô} — the \emph{m} of English \emph{met}, plus the inherent vowel.

Three consonants now. Say them in the order you met them, each with the vowel nobody wrote:

\begin{quote}
\textbf{\bn{ন}}  \textbf{\bn{হ}}  \textbf{\bn{ম}}
\end{quote}

\emph{nô, hô, mô.} Notice you are saying \emph{aw} three times, not \emph{ah}.

\subsection*{Your turn}
\begin{itemize}
\raggedright
  \item \emph{Write it:} \bn{ম}
  \item \emph{Read it:} every piece this chapter has given you so far, in order
\end{itemize}

\begin{tcolorbox}[breakable,colback=teal!4,colframe=teal!35!black,title={Before you move on}]
Say the three consonants you hold, each with its inherent vowel. (\emph{\textbf{nô, hô, mô}}.)

Source: \href{https://www.unicode.org/charts/PDF/U0980.pdf}{Unicode Bengali chart}.
\end{tcolorbox}
\section[shô]{\bn{স} — the consonant that is written s and said sh}

\label{lesson:BN-W01-sa}



\begin{quote}
Write the piece from the previous lesson before adding another.
\end{quote}

\subsection*{Script}
\begin{quote}
\bn{স}
\end{quote}

Here the letter and the sound part company, and it is worth knowing early.

This letter descends from the Sanskrit \textbf{s}, and every transliteration writes it \emph{s}. \textbf{In Bengali it is normally pronounced \emph{sh}} — the \emph{sh} of English \emph{shoe}.

So the word you say as \emph{nomosh-kar} is written with what looks like an \emph{s}. The spelling records the word's Sanskrit ancestry; the pronunciation records what Bengali did with it afterwards. \textbf{Both are true, and the script keeps the older one.}

English does exactly this with \emph{knight} and \emph{through}. Every literary language accumulates spellings that outlive their sounds.

\subsection*{Your turn}
\begin{itemize}
\raggedright
  \item \emph{Write it:} \bn{স}
  \item \emph{Read it:} every piece this chapter has given you so far, in order
\end{itemize}

\begin{tcolorbox}[breakable,colback=teal!4,colframe=teal!35!black,title={Before you move on}]
How is \textbf{\bn{স}} written in transliteration, and how is it said? (Written \emph{\textbf{s}}, said \emph{\textbf{sh}}.) Why the difference? (The spelling keeps the \textbf{Sanskrit ancestry}; the sound moved on.)

Source: \href{https://www.unicode.org/charts/PDF/U0980.pdf}{Unicode Bengali chart}.
\end{tcolorbox}
\section[kô]{\bn{ক} — the consonant ka}

\label{lesson:BN-W01-ka}



\begin{quote}
Write the piece from the previous lesson before adding another.
\end{quote}

\subsection*{Script}
\begin{quote}
\bn{ক}
\end{quote}

\textbf{kô} — the \emph{k} of English \emph{skip}, with \textbf{no puff of breath} after it.

Hold a palm in front of your mouth and say English \emph{kit}: you feel air. Say \emph{skip}: you do not. \textbf{This is the second one.} Bengali has a separate letter for the breathy version, and swapping them changes the word — so a difference English treats as an accident, this script writes down.

Five consonants, one mātrā, one independent vowel. One piece of machinery left.

\subsection*{Your turn}
\begin{itemize}
\raggedright
  \item \emph{Write it:} \bn{ক}
  \item \emph{Read it:} every piece this chapter has given you so far, in order
\end{itemize}

\begin{tcolorbox}[breakable,colback=teal!4,colframe=teal!35!black,title={Before you move on}]
Is \textbf{\bn{ক}} the \emph{k} of \emph{kit} or of \emph{skip}? (Of \emph{\textbf{skip}} — no puff of air.) Does the script care? (\textbf{Yes} — the breathy one is a different letter.)

Source: \href{https://www.unicode.org/charts/PDF/U0980.pdf}{Unicode Bengali chart}.
\end{tcolorbox}
\section[(hasanta)]{\bn{্} — the mark that DELETES the inherent vowel and fuses two consonants}

\label{lesson:BN-W01-hasanta}



\begin{quote}
Write the piece from the previous lesson before adding another.
\end{quote}

\subsection*{Script}
The inherent vowel is a nuisance when you do not want it. If a word presses two consonants together with no vowel between them, writing them plainly would insert an \textbf{ô} that is not there.

So there is a mark that removes it:

\begin{quote}
\bn{্}
\end{quote}

Bengali calls it the \textbf{hasanta} — Devanagari calls the same thing a \emph{virama}. It sits under a consonant and means: \emph{no vowel here.}

Then something happens that no European alphabet does. A consonant stripped of its vowel \textbf{fuses} with the one after it into a single joined shape, a \textbf{conjunct}:

\begin{quote}
\textbf{\bn{স}} + \textbf{\bn{্}} + \textbf{\bn{ক}} $\to$ \textbf{\bn{স্ক}}
\end{quote}

One cluster. \textbf{Conjuncts are not new characters to memorise} — they are two you already know, joined. Bengali's are more thoroughly fused than Devanagari's, so some take a moment to take apart, but the principle does not change.

\subsection*{Your turn}
\begin{itemize}
\raggedright
  \item \emph{Write it:} \bn{্}
  \item \emph{Read it:} every piece this chapter has given you so far, in order
\end{itemize}

\begin{tcolorbox}[breakable,colback=teal!4,colframe=teal!35!black,title={Before you move on}]
What does the hasanta do? (\textbf{Removes the inherent vowel}.) What is Devanagari's name for the same mark? (The \textbf{virama}.) And are conjuncts new letters? (\textbf{No} — two known letters, fused.)

Source: \href{https://www.unicode.org/charts/PDF/U0980.pdf}{Unicode Bengali chart}.
\end{tcolorbox}
\section[rô]{\bn{র} — the last consonant of the greeting}

\label{lesson:BN-W01-ra}



\begin{quote}
Write the piece from the previous lesson before adding another.
\end{quote}

\subsection*{Script}
\begin{quote}
\bn{র}
\end{quote}

\textbf{rô} — a light tap of the tongue, closer to the \emph{r} of Spanish \emph{pero} than to any English \emph{r}.

That is the last piece. Everything the greeting needs is now something you can write from memory.

\subsection*{Your turn}
\begin{itemize}
\raggedright
  \item \emph{Write it:} \bn{র}
  \item \emph{Read it:} every piece this chapter has given you so far, in order
\end{itemize}

\begin{tcolorbox}[breakable,colback=teal!4,colframe=teal!35!black,title={Before you move on}]
How is \textbf{\bn{র}} made? (A \textbf{light tap}, not an English \emph{r}.)

Source: \href{https://www.unicode.org/charts/PDF/U0980.pdf}{Unicode Bengali chart}.
\end{tcolorbox}
\section[nomoshkar]{\bn{নমস্কার} — seven pieces, all of them already yours}

\label{lesson:BN-W01-nomoshkar-read}



\begin{quote}
Write the piece from the previous lesson before adding another.
\end{quote}

\subsection*{Script}
Nothing new is introduced here. Every piece below is one you can write.

\begin{quote}
\textbf{\bn{ন}} + \textbf{\bn{ম}} + \textbf{\bn{স}} + \textbf{\bn{্}} + \textbf{\bn{ক}} + \textbf{\bn{া}} + \textbf{\bn{র}} $\to$ \textbf{\bn{নমস্কার}}
\end{quote}

Read it the way you built it:

\begin{itemize}
\raggedright
  \item \textbf{\bn{ন}} — \emph{nô}, inherent vowel
  \item \textbf{\bn{ম}} — \emph{mô}, inherent vowel
  \item \textbf{\bn{স্ক}} — the conjunct: the \emph{s}-letter with its vowel removed, fused to \emph{kô}
  \item \textbf{\bn{া}} — the mātrā, making that piece \emph{ka}
  \item \textbf{\bn{র}} — \emph{rô}
\end{itemize}

\textbf{nô-mosh-kar} — and notice the \emph{sh}, from a letter written \emph{s}.

Count what it took: \textbf{six consonants, one mātrā, one independent vowel and one hasanta.} With them you can read \emph{no} and the greeting, and every consonant you meet from here obeys rules you already hold.

Bengali has around forty more letters. It does not have any more \textbf{systems}.

\subsection*{Your turn}
\begin{itemize}
\raggedright
  \item \emph{Write it:} \bn{নমস্কার}
  \item \emph{Read it:} every piece this chapter has given you so far, in order
\end{itemize}

\begin{tcolorbox}[breakable,colback=teal!4,colframe=teal!35!black,title={Before you move on}]
How many new pieces were in this lesson? (\textbf{None}.) What is \textbf{\bn{স্ক}} made of? (The \textbf{s}-letter with its vowel removed, joined to \textbf{\bn{ক}}.) And what still remains to learn — rules, or shapes? (\textbf{Shapes}.)

Source: \href{https://www.unicode.org/charts/PDF/U0980.pdf}{Unicode Bengali chart}.
\end{tcolorbox}
